% s-125-p125-m001; original PDF 125
\Aphorism{s-125-p125-m001}{162}
Die meisten Ehen missrathen, weil sie von Anfang an auf eine kleinliche Nähe und Intimität gestellt sind: „Alles mit einander gemein haben“ ist ein gefahrvolles Wagniss, bei dem bald aus der Gemeinschaft eine Gemeinheit wird. Man soll die zarten Farben der Ferne, den Duft und Schmelz der ersten blöden Verliebtheit, die Discretion des vorehelichen Verkehrs in die Ehe hinüberretten, gerade weil es hier eine letzte Vertraulichkeit giebt, die zur Quelle tiefster wechselseitiger Abneigung werden kann. Müssen sich Eheleute durchaus täglich, stündlich, allaugenblicklich sehen, muss man einander zu Bewusstsein bringen, dass der Mensch ein absonderndes, ausdünstendes und zu Zeiten unappetitliches Thier ist? Wenn Herr und Dame sich in Mann und Weib verwandeln, so mögen sie zugleich an Stelle der gesellschaftlichen gêne ihre ganze erfinderische Scham aufbieten, um vor einander möglich zu bleiben.
\par

